\chapter{Tratamientos térmicos y control microestructural}\label{ch:08}
\objetivos{Explicar qué cambia durante un tratamiento y por qué su secuencia importa; distinguir endurecimiento por transformación, precipitación y deformación.}
\section{El tratamiento como historia}
Un tratamiento térmico es una historia de temperatura y tiempo, con condiciones de calentamiento, permanencia, enfriamiento y atmósfera. No se define solo por su temperatura máxima. La sección de la pieza, la transferencia térmica y la carga del horno condicionan gradientes y velocidad de enfriamiento; por eso una receta de probeta no se traslada sin más a un componente masivo.

La microestructura inicial también importa. Dos piezas de la misma composición nominal pueden responder de forma distinta si tienen deformación previa, segregación, tamaño de grano o distribución de fases diferentes. El control incluye trazabilidad del lote, equipo, instrumentación y verificación de propiedades.
\section{Aceros: temple y revenido}
En una secuencia convencional para aceros susceptibles de endurecimiento se austenitiza y se enfría con rapidez suficiente para obtener una estructura martensítica en la zona de interés. La templabilidad describe la capacidad de endurecer en profundidad bajo condiciones de enfriamiento, y no es sinónimo de dureza máxima. Composición, sección y medio de temple intervienen.

El revenido modifica la martensita y los precipitados, alivia parte de las tensiones y ajusta dureza, resistencia y tenacidad. No es una operación decorativa posterior al temple. La combinación final depende de temperatura y tiempo; existen intervalos de fragilización y comportamientos específicos que impiden generalizar que «más revenido siempre mejora todo».

Un enfriamiento desigual puede introducir tensiones residuales, distorsión y grietas. Esas tensiones se suman a las de servicio en el análisis adecuado. Un resultado de dureza satisfactorio en una superficie no demuestra por sí solo uniformidad de propiedades en toda la sección.
\section{Solubilización, temple y precipitación}
En una aleación endurecible por precipitación se busca disolver determinadas fases a una temperatura de solubilización adecuada, enfriar para retener una solución sobresaturada y madurar para desarrollar una distribución de precipitados. El temple de esta secuencia tiene la función de retener soluto; no debe confundirse con la formación de martensita de un acero.

La maduración natural ocurre a temperatura ambiente en sistemas que la presentan; la artificial utiliza una temperatura elevada controlada. La resistencia puede aumentar hasta un máximo y descender con sobreenvejecimiento por evolución del tamaño, coherencia y separación de precipitados. Acelerar el proceso es solo parte de la explicación: el estado alcanzado y sus compromisos también cambian.
\ilustra{tratamiento-termico}{Historia esquemática de solubilización, enfriamiento y maduración, junto a una evolución cualitativa de resistencia. No constituye una receta industrial.}
\section{Tiempo y temperatura no se intercambian libremente}
La difusión y la nucleación compiten con crecimiento y transformaciones. Una temperatura mayor puede acelerar procesos y cambiar las fases presentes. Las correlaciones de tiempo--temperatura se utilizan dentro de dominios calibrados, no como una equivalencia universal entre cualquier tratamiento.

Los diagramas de fases describen estados de equilibrio; los diagramas TTT o CCT incorporan cinética bajo historias concretas. Leer una fracción de equilibrio no demuestra que esa fracción se haya alcanzado en un enfriamiento rápido. El material real puede conservar fases metaestables o gradientes por difusión incompleta.
\section{Procesamiento termomecánico y superficie}
La deformación seguida de recuperación y recristalización modifica densidad de dislocaciones, textura y tamaño de grano. La combinación de forja y tratamiento controla propiedades y defectos de una manera que no se resume solo con la temperatura final. La orientación del flujo de forja puede afectar la respuesta en direcciones distintas.

Los tratamientos superficiales modifican una capa y pueden buscar dureza, resistencia al desgaste, protección ambiental o tensiones residuales favorables. Cementación y nitruración alteran química superficial; granallado introduce deformación y tensiones de compresión; los recubrimientos añaden otra fase e interfaces. Cada proceso requiere comprobar compatibilidad con geometría, defectos y condiciones térmicas posteriores.
\caso{M08. Una reparación puede alterar el estado}{Una zona de una pieza ficticia endurecida por precipitación se calienta localmente durante una intervención. Aunque no cambie su composición nominal, pueden disolverse, crecer o transformarse precipitados. Medir solo dimensiones no demuestra que se conserven las propiedades previas. La evaluación necesita la historia térmica, el material, el estado inicial y los datos de reparación aplicables; el libro no autoriza un tratamiento de recuperación.}
\section{Qué debe documentar un problema}
Un ejercicio de tratamientos debe especificar aleación, estado inicial y objetivo. Si se presenta un diagrama simplificado, debe indicar si es de equilibrio o cinético y qué condiciones representa. Una comparación puede explicar tendencias sin proponer una receta de horno para una pieza aeronáutica.

En fabricación y mantenimiento, la conformidad se apoya en procesos aprobados y registros, no en reproducir de memoria una temperatura de un manual. En el estudio, el objetivo consiste en predecir qué variables cambian, formular hipótesis y reconocer cuándo una medición es insuficiente.
\section{Tratamiento térmico de aleaciones de fundición}
La afirmación de que los aluminios moldeados no admiten tratamientos térmicos es incorrecta. Existen aleaciones de aluminio de fundición, incluidas determinadas Al--Si--Mg, que pueden recibir solubilización, temple y envejecimiento artificial y suministrarse en estados como T6. Su aplicación estructural debe justificarse de forma separada mediante composición, defectos de colada, control de proceso y especificación aplicable.

En una aleación endurecible, el enfriamiento desde solubilización sirve para retener soluto; no equivale al temple martensítico de un acero. La maduración artificial no solo acelera la natural, también modifica el estado de precipitación y el equilibrio entre resistencia, tenacidad y comportamiento ambiental. Una excursión térmica de servicio puede alterar propiedades sin deformación visible; una pieza aparentemente intacta no demuestra conservación de su condición original.


\section{Un diagrama de equilibrio no es una historia de proceso}
Un diagrama de fases representa equilibrios para composición y temperatura bajo las condiciones del modelo. Una línea de enlace permite identificar composiciones de fases coexistentes; la regla de la palanca proporciona sus fracciones en masa cuando el eje de composición usa esa base. Para una aleación global $C_0$ entre $C_\alpha$ y $C_\beta$,
\begin{equation}w_\alpha=\frac{C_\beta-C_0}{C_\beta-C_\alpha},\qquad w_\beta=\frac{C_0-C_\alpha}{C_\beta-C_\alpha}.\end{equation}
Ambas fracciones suman uno y permanecen entre cero y uno dentro de la región bifásica. No son automáticamente fracciones volumétricas: convertirlas exige densidades. Fuera de esa región no se aplica la misma construcción.

Un enfriamiento rápido puede impedir que se alcance el equilibrio. Por eso un diagrama por sí solo no da dureza, tamaño de precipitado o fracción transformada tras un tiempo concreto. La cinética requiere información adicional; las curvas tiempo--temperatura--transformación y las de enfriamiento continuo responden a preguntas diferentes. No se emplean como calendarios intercambiables.
\ilustra{diagrama-fases}{Diagrama binario isomorfo esquemático con región líquido más sólido. Una línea de enlace en la región bifásica permite aplicar la regla de la palanca. No representa un sistema real ni permite definir un tratamiento industrial.}
\section{Precipitación: nucleación, crecimiento y sobreenvejecimiento}
Una secuencia típica de endurecimiento por precipitación parte de una solución sólida sobresaturada que evoluciona durante maduración. Precipitados finos y su relación con la matriz dificultan movimiento de dislocaciones. Al crecer y cambiar coherencia y separación puede variar el mecanismo de interacción y caer la resistencia tras un máximo. La explicación exige identificar la aleación: no todas las secuencias ni fases son iguales.

Una mayor dureza no demuestra automáticamente mejor tenacidad, corrosión bajo tensión o estabilidad térmica. El tiempo de servicio también puede actuar como exposición térmica y modificar el estado obtenido en fábrica. Si una reparación aplica calor local, el efecto puede extenderse más allá de la zona visible; no se restituye una condición solo por observar el mismo color de superficie.
\section{Control del proceso y de sus resultados}
Un proceso reproducible define identificación de pieza, estado inicial, equipos, medio, temperatura, permanencia, transferencia y enfriamiento, además de inspecciones posteriores. Las tolerancias y verificaciones proceden de una especificación aplicable; no se inventan a partir de un esquema docente. La evidencia combina registro del proceso y resultados pertinentes: dimensiones, dureza, microestructura u otros ensayos según la exigencia.
\repaso{Si se aumenta mucho el tiempo de maduración, ¿la resistencia debe seguir creciendo?}{No. Puede alcanzarse un máximo y después sobreenvejecimiento por evolución de precipitados. La respuesta depende de aleación, temperatura y estado previo; la dureza medida debe relacionarse con el requisito específico.}

\resumen{Temple no significa lo mismo en todas las familias. La solubilización y la maduración forman una secuencia. Sobreenvejecer no equivale simplemente a «tratar peor»: puede perseguir otro compromiso, pero requiere datos. Una intervención térmica puede alterar propiedades aunque no deje una deformación visible.}
\fuentes{FAA--General, cap. 7, tratamientos de metales ferrosos y no ferrosos \cite{FAA-G}; materiales docentes de introducción; mecanismos de dislocaciones y recuperación \cite{MIT-CREEP}.}
